
% SGA 3，报告 VII_B；英文重建本第一轮的中文译本。
% 范围： 节 2.8 全文 依据文献局部页 63--65,
% 止于下节之前： 节 2.9.
% 合并阅读本页： 576--578. 重版印刷页： 565--567.
% 范围内编者注： 115--117.
% 依据： Polo--Gille Exp7B-13oct24.pdf,
% SHA-256 D303690E0195654BBE7C3DAA5D89EA23E24CD709A3575217234AC19AE8836694.
% 参照了 Jacob Reinhold 的 CC BY 4.0 Markdown，作为对照来源。
% 本组件无图；所有数学内容均为原生 LaTeX。
\subsection*{2.8. 域上的幺幂群概形}
\phantomsection
\label{sga3:VIIB:2.8}
\addcontentsline{toc}{subsection}{\SGARefLink{sga3:VIIB:2.8}{2.8}. 域上的幺幂群概形}
\footnote{\label{sga3:VIIB:nde:115}此小节由编者添加。}
设 \(k\) 是域。回顾，若 \(k\)-群概形 \(G\) 满足以下两个条件，则称其为\emph{幺幂}的\textup{（参见 [DG70]，\S~IV.2，\SGARefLink{sga3:VIIB:2.5:proposition}{命题 2.5}）}：
\begin{enumerate}
\item[\textup{(a)}] \(G\) 是仿射的。
\item[\textup{(b)}] 每个单 \(\mathcal O(G)\)-余模都是平凡的。具体地，若
\[
  \rho:V\longrightarrow V\otimes_k\mathcal O(G)
\]
是非零 \(k\)-向量空间 \(V\) 上的 \(\mathcal O(G)\)-余模结构，且不存在使
\[
  \rho(W)\subset W\otimes_k\mathcal O(G),
\]
成立的非零真子空间 \(W\subset V\)，则 \(V\) 是一维的，且
\[
  \rho(v)=v\otimes1
  \qquad\text{对每个 }v\in V.
\]
\end{enumerate}
根据所引结果，当 \(G\) 在 \(k\) 上有限型时，这等价于报告 XVII，\S~1 中给出的定义：\(G\) 有一个有限合成列，其相继商同构于 \(\mathbb G_{a,k}\) 的 \(k\)-子群。
对每个仿射 \(k\)-群概形 \(G\)，\(\mathcal O(G)\) 的余乘法赋予伪紧 \(k\)-模
\[
  A=\mathcal O(G)^*
\]
一个不一定交换的预有限 \(k\)-代数结构；其单位元 \(1_A\) 就是增广
\[
  \varepsilon:\mathcal O(G)\longrightarrow k.
\]
置
\[
  I=\{f\in A\mid f(1_{\mathcal O(G)})=0\}.
\]
则 \(I\) 是 \(A\) 的双边理想，且
\[
  A=k\cdot1_A\oplus I;
\]
参见\ \SGARefLink{sga3:VIIB:1.3.6}{1.3.6}。
设 \(V\) 是 \(A\) 的有限余维子空间，考察连续 \(k\)-双线性映射
\[
  \varphi:A\times A\longrightarrow A/V,
  \qquad
  (a,b)\longmapsto ab+V.
\]
由\SGARefLink{sga3:VIIB:0.3.1}{引理 0.3.1}，存在有限余维子空间 \(L_1,L_2\subset A\)，使 \(V\) 包含 \(AL_2\) 和 \(L_1A\)。于是
\[
  L=L_1\cap L_2
\]
是有限余维的，且 \(V\) 包含双边理想 \(ALA\)，后者本身也是有限余维的。因此，有限余维双边理想构成 \(A\) 中零点的一个邻域基。
由此可得，\(\mathcal O(G)\)-余模与连续 \(A\)-模是一回事：后者是一个赋予离散拓扑的 \(A\)-模 \(V\)，使得映射
\[
  A\times V\longrightarrow V
\]
连续。每个这样的模都是有限维 \(A\)-子模 \(V_i\) 的并，而每个 \(V_i\) 都是 \(A\) 的某个有限维商 \(k\)-代数 \(A_i\) 上的模。
因此，单连续 \(A\)-模 \(M\) 在 \(k\) 上有限维，并且是有限维代数 \(A/\operatorname{Ann}(M)\) 上的忠实单模。因此，后者为有限维单 \(k\)-代数，故 \(\operatorname{Ann}(M)\) 是 \(A\) 的开极大理想。
反之，设 \(P\) 是 \(A\) 的开素理想。\footnote{\label{sga3:VIIB:nde:116}回顾，若在 \(A/P\) 中任意两个非零双边理想的积都非零，就称环 \(A\) 的双边理想 \(P\) 为素理想。}
则 \(A/P\) 是零理想为素理想的有限维 \(k\)-代数，因此是有限维单 \(k\)-代数。故在同构意义下，恰好存在唯一一个零化子为 \(P\) 的单连续 \(A\)-模。于是
\[
  M\longmapsto\operatorname{Ann}(M)
\]
是从单连续 \(A\)-模的同构类到 \(A\) 的开素理想的双射。
特别地，在 \(k\) 上一维的 \(A\)-模 \(A/I\) 称为\emph{平凡模}。它对应于平凡的一维 \(\mathcal O(G)\)-余模 \(V\)，对后者有
\[
  \rho(v)=v\otimes1_{\mathcal O(G)}
  \qquad\text{对每个 }v\in V.
\]
我们证明了：
\medskip
\noindent\textbf{命题 2.8.1。}
\phantomsection
\label{sga3:VIIB:2.8.1}
{\itshape 设 \(k\) 是域，\(G\) 是仿射 \(k\)-群概形，置 \(A=\mathcal O(G)^*\)。
\begin{enumerate}
\item[\textup{(i)}] \(G\) 是幺幂的当且仅当 \(I\) 是 \(A\) 的唯一开素理想。\footnote{\label{sga3:VIIB:nde:117}等价地，\(k\cdot1_{\mathcal O(G)}\) 是 \(\mathcal O(G)\) 的唯一单 \(k\)-子余代数；例如参见 [Ab80]，3.1.4。编者还指出所引文献第 130 页第 4 行的一个误印：应将 \(M\simeq C^*/\operatorname{ann}M\) 改为 \(C^*/\operatorname{ann}M\simeq\operatorname{End}_k(M)\)。}
\item[\textup{(ii)}] 特别地，设 \(G\) 是交换的，从而
\[
  \mathcal O(G)=\mathbf H(D(G)),
\]
其中 \(D(G)=\operatorname{Spf}(A)\) 是 \(G\) 的 Cartier 对偶。则 \(G\) 是幺幂的当且仅当 \(D(G)\) 是无穷小的。
\end{enumerate}}
